\chapter{Stakeholder Engagement}
\label{ap:stakeholder}

\newtcolorbox{emailbox}[1]{
    colback=white,
    colframe=gray!30,
    fonttitle=\bfseries\sffamily,
    coltitle=black,
    title={#1},
    enhanced,
    attach boxed title to top left={yshift=-2mm, xshift=2mm},
    boxed title style={colback=gray!10},
    breakable
}

Communication with \gls{fps} primarily occurred over email between:
\begin{itemize}
    \item Diego Temkin, undergraduate student, Department of Electrical Engineering and Computer Science and Department of Urban Studies and Planning
    \item Riccardo Fiorista, PhD student in Transportation, Department of Civil and Environmental Engineering and Laboratory for Information and Decision Systems
    \item John Attanuchi, Research Associate, Center for Transportation and Logistics
    \item Lorena Tovar, Assistant Director of Academic Programs Administration in Department of Urban Studies and Planning, Framingham School Committee member from District 2
    \item Dr. Robert Tremblay, Superintendent of Schools at \gls{fps}
    \item Lincoln Lynch, IV, Executive Director of Finance and Operations at \gls{fps}
\end{itemize}

\section{Initial Email Outreach}

The following emails occurred between November 19th, 2025, and December 27th, 2025.

{
\small
\begin{emailbox}{Framingham School bus service improvement}
    \textbf{From:} Riccardo Fiorista \\
    \textbf{To:} Lorena Tovar \\
    \textbf{Cc:} John P Attanuchi; Jinhua Zhao; Alexandre Jacquillat; Diego Temkin \\
    \textbf{Subject:} Framingham School bus service improvement

    \vspace{1em}
    Hi Lorena, \vspace{1em}

    It was good running into you yesterday, and congratulations again on your election. As we discussed, some political groundwork may need to come first, but we are ready to support the effort with our tools. We have already pulled the publicly available data (school locations, the street network, public transit schedules, and census-block demographics) and begun sketching a possible problem formulation. \vspace{1em}

    To move toward an actual solution, we will need more specific information: where the students who rely on school bus service live (street-level detail would be great if addresses are not available), and whether any of them fall into vulnerable groups. It would also help to know what planning tools the district currently uses and to see the existing routing. \vspace{1em}

    John suggested that we come to Framingham to meet with you and the relevant stakeholders. That would give us the context we need and allow us to shape an algorithmic solution that fits the district’s constraints. \vspace{1em}

    Let us know if a meeting this year is still feasible. \vspace{1em}
    
    Best regards, \\
    Riccardo
\end{emailbox}
}

\section{Stakeholder Meeting}

On January 8th, 2026, the group from MIT including Diego Temkin, John Attanuchi, and Riccardo Fiorista (virtual) met with Dr. Robert Tremblay, Lincoln Lynch, and Lorena Tovar at the \gls{fps} administrative offices. Handwritten notes were taken, and the following primary questions were answered:

\begin{itemize}
    \item Are students zoned to a specific school or does Framingham have open enrollment? School choice is prevalent in Framingham. Students attending a school over 2 miles away are given bus services for free, and those below can sign up. 
    \item Are school buses centrally managed or does every school have it's [sic] own fleet? School buses are centrally managed by the district, but typically only go to one school destination.
    \item How do people sign up for school buses, and how are they assigned a route? For the next academic year, all students are automatically signed up for school bus service, but are able to opt out.
    \item How are school bus routes created currently? Would we be able to get the current plans? They have a platform that allows them to create routes, but they are manually adjusted if new students sign up.
    \item How do students with disabilities use school buses in Framingham? Are there any special considerations or constraints? Special monitors are available for students with disabilities, and those students are encouraged to ride in a normal school bus.
    \item Is there any integration with other transit providers, such as \gls{mwrta}? No, and \gls{mwrta} routes are not convenient for getting to many schools in Framingham.
\end{itemize}

\section{Further Information Requests}

{
\small
\begin{emailbox}{Thank you + Next Steps}
    \textbf{From:} Lorena Tovar \\
    \textbf{To:} John P Attanuchi; Riccardo Fiorista; Diego Temkin; Robert Tremblay; Lincoln Lynch  \\
    \textbf{Subject:} Thank you + Next Steps

    \vspace{1em}
    Hi All, \vspace{1em}

    Thanks again for making this meeting possible-\/- especially John and Diego for coming out to Framingham, and Riccardo for joining us from Vienna. \vspace{1em}
    
    Now that we are all connected, here are the next steps: \vspace{1em}
    
    1. \acrshort{fps} will connect with the vendor and share the \acrshort{gps} tracking info, if possible. \\
    2. \acrshort{fps} will share all relevant data \\
    3. MIT will work with the data in this `exploratory data analysis phase', then update the group on whether it's necessary to reconvene or to outline the next steps. \vspace{1em}
    
    Bob and Lincoln-\/- is there a way that we can look at truancy and tardiness data to see if there's a correlation between truancy and bus access? This wasn't in the original scope we discussed yesterday, but it occurred to me that we do have an issue with truancy/tardiness so it might be worth exploring. \vspace{1em}
    
    Thanks again! \\
    Lorena
\end{emailbox}

\begin{emailbox}{Thank you + Next Steps}
    \textbf{From:} Lincoln Lynch \\
    \textbf{To:} Lorena Tovar \\
    \textbf{Cc:} John P Attanuchi; Riccardo Fiorista; Diego Temkin; Robert Tremblay \\
    \textbf{Subject:} Re: Thank you + Next Steps

    \vspace{1em}
    Hi Lorena. Thank you all for meeting with us yesterday. I do not have truancy or tardiness data but we can talk to our truant officer to see if we can obtain data. \vspace{1em}

    I am working on obtaining a student privacy agreement that I will send over. As soon as that is completed I can get the student data over. Have a great weekend!
    
    \noindent\rule{\textwidth}{1pt}

    Lincoln Lynch IV, Executive Director of Finance and Operations \\
    \Acrlong{fps} - Business Office \\
    \textit{[district office address redacted]} \\
    \textit{[email redacted]} \vspace{1em}

    \Acrlong{fps} provides equal employment opportunities (EEO) to all employees and applicants for employment without regard to race, color, religion, sex, national origin, age, disability or genetics. In addition to federal law requirements, \acrlong{fps} complies with applicable state and local laws governing nondiscrimination in employment in every location in which the District has facilities. This policy applies to all terms and conditions of employment, including recruiting, hiring, placement, promotion, termination, layoff, recall, transfer, leaves of absence, compensation and training. \acrlong{fps} expressly prohibits any form of workplace harassment based on race, color, religion, gender, sexual orientation, gender identity or expression, national origin, age, genetic information, disability, or veteran status.
\end{emailbox}

\begin{emailbox}{Thank you + Next Steps}
    \textbf{From:} John P Attanuchi \\
    \textbf{To:} Lincoln Lynch; Lorena Tovar \\
    \textbf{Cc:} Riccardo Fiorista; Diego Temkin; Robert Tremblay \\
    \textbf{Subject:} Re: Thank you + Next Steps

    \vspace{1em}
    Thanks Lincoln and Lorena for a great meeting. You can send the non-disclosure agreement to me to sign on behalf of the MIT research team. We look forward to getting started with the data as soon as possible this month as the students don't start their Spring semester classes until the beginning of February and they have a lot of time for research. Regarding the bus \acrshort{gps} data, I had a thought that, if it would be easier for you, \textit{[paraphrased: John offered alternative options for transferring the bus GPS data, including direct system access if convenient.]}. Just a thought but no problem if it can't be done. Thanks again.
    
    John \vspace{1em}
    
    \textit{[mobile email signature redacted]}
\end{emailbox}

\begin{emailbox}{Thank you + Next Steps}
    \textbf{From:} Lincoln Lynch \\
    \textbf{To:} Lorena Tovar \\
    \textbf{Cc:} John P Attanuchi; Riccardo Fiorista; Diego Temkin; Robert Tremblay \\
    \textbf{Subject:} Re: Thank you + Next Steps

    \vspace{1em}
    Hi, good morning John. I am working on getting the student list and have pretty much solidified it. We are working on \lowercase{\acrshort{gps}} data as well. We do not give non-\acrlong{fps} access to our systems so I do not think we would do that. \vspace{1em}

    After talking internally, we usually obtain a data privacy agreement from the requester of data. Does MIT have a standard agreement that basically says you will not use the data for any purposes other than the research?
    
    \noindent\rule{\textwidth}{1pt}

    Lincoln Lynch IV, Executive Director of Finance and Operations \\
    \Acrlong{fps} - Business Office \\
    \textit{[district office address redacted]} \\
    \textit{[email redacted]} \vspace{1em}

    \textit{[\gls{fps} standard EEO footer; reproduced once here, omitted from subsequent emails.]}
\end{emailbox}

\begin{emailbox}{Thank you + Next Steps}
    \textbf{From:} John P Attanuchi \\
    \textbf{To:} Lincoln Lynch \\
    \textbf{Cc:} Lorena Tovar; Riccardo Fiorista; Diego Temkin; Robert Tremblay \\
    \textbf{Subject:} Re: Thank you + Next Steps \\
    \textbf{Attached:} \textit{[attachment name redacted]}

    \vspace{1em}
    Hi Lincoln,

    I should have thought of this sooner! Attached please find the signed standard MIT research sponsor confidential information agreement. I don’t have the authority to change or negotiate the specific terms of this agreement, but it looks like it should protect your data from any disclosure or distribution by anyone here at MIT. I hope that you agree and we can receive your data soon to begin our exploratory research. Thanks again.

    John
\end{emailbox}

\begin{emailbox}{Thank you + Next Steps}
    \textbf{From:} Lincoln Lynch \\
    \textbf{To:} John P Attanuchi \\
    \textbf{Cc:} Lorena Tovar; Riccardo Fiorista; Diego Temkin; Robert Tremblay \\
    \textbf{Subject:} Re: Thank you + Next Steps \\
    \textbf{Attached:} \textit{[attachment name redacted]}

    \vspace{1em}
    Hi John. Thanks for this! I signed off on it. Also attached is the requested student-level dataset to start. I am working on the bus data as well.
    
    \noindent\rule{\textwidth}{1pt}

    Lincoln Lynch IV, Executive Director of Finance and Operations \\
    \Acrlong{fps} - Business Office \\
    \textit{[district office address redacted]} \\
    \textit{[email redacted]} \vspace{1em}

    \textit{[\gls{fps} standard EEO footer; reproduced once here, omitted from subsequent emails.]}
\end{emailbox}

\begin{emailbox}{Thank you + Next Steps}
    \textbf{From:} John P Attanuchi \\
    \textbf{To:} Lincoln Lynch \\
    \textbf{Cc:} Lorena Tovar; Riccardo Fiorista; Diego Temkin; Robert Tremblay \\
    \textbf{Subject:} Re: Thank you + Next Steps \\
    \textbf{Attached:} \textit{[attachment name redacted]}

    \vspace{1em}
    Wonderful-\/-Diego and Riccardo will start looking at the data and get back to you if they have any questions. One more request in the interim may be a list of all the stops on the current bus routes in advance of the \acrshort{gps} data if that is available in any form. Thanks again.

    John \vspace{1em}

    \textit{[mobile email signature redacted]}
\end{emailbox}

\begin{emailbox}{Thank you + Next Steps}
    \textbf{From:} Riccardo Fiorista \\
    \textbf{To:} Lincoln Lynch \\
    \textbf{Cc:} Lorena Tovar; Diego Temkin; Robert Tremblay; John P Attanuchi \\
    \textbf{Subject:} Re: Thank you + Next Steps \\
    \textbf{Attached:} \textit{[attachment name redacted]}

    \vspace{1em}
    Hi Lincoln, \vspace{1em}

    Thanks for sending through the data! \vspace{1em}

    We ran a first pass on the file and were able to geocode all student addresses. One thing that jumped out is that we’re seeing roughly 144 addresses outside of Framingham, including some as far as Worcester and even Cape Cod. Can you confirm if that’s expected, or if we should treat some of those as data issues (old addresses, placements, typos, etc.)? \vspace{1em}

    Also, we weren’t able to fully map all school codes to precise school names and locations. \\
    Could you please help us to revise the attached school list and:

    \begin{itemize}
        \item Confirm we identified the correct school names/addresses for the codes we did map, and
        \item Fill in the missing ones (we also noticed at least one entry that looks like Framingham History Center, which seems wrong).
    \end{itemize}

    One optional request and only if it’s feasible: Do you have any indicator for students who require a bus monitor? If that’s not readily available, no worries at all. \vspace{1em}

    Thanks and have a great weekend, \\
    Riccardo
\end{emailbox}

\begin{emailbox}{Thank you + Next Steps}
    \textbf{From:} Lincoln Lynch \\
    \textbf{To:} Riccardo Fiorista \\
    \textbf{Cc:} Lorena Tovar; Diego Temkin; Robert Tremblay; John P Attanuchi \\
    \textbf{Subject:} Re: Thank you + Next Steps \\
    \textbf{Attached:} \textit{[attachment name redacted]}

    \vspace{1em}
    Hi Ricardo [sic]. Hope all is well. Attached is an excel with a column I added with my notes. \textit{[paraphrased: Lincoln explained that out-of-district addresses correspond to \gls{mkv} students who are transported separately by minivan and can be excluded from the school-bus analysis.]}
    
    \noindent\rule{\textwidth}{1pt}

    Lincoln Lynch IV, Executive Director of Finance and Operations \\
    \Acrlong{fps} - Business Office \\
    \textit{[district office address redacted]} \\
    \textit{[email redacted]} \vspace{1em}

    \textit{[\gls{fps} standard EEO footer; reproduced once here, omitted from subsequent emails.]}
\end{emailbox}

\begin{emailbox}{Thank you + Next Steps}
    \textbf{From:} Diego Temkin \\
    \textbf{To:} Lincoln Lynch; Riccardo Fiorista \\
    \textbf{Cc:} Lorena Tovar; Robert Tremblay; John P Attanuchi \\
    \textbf{Subject:} Re: Thank you + Next Steps

    \vspace{1em}
    Hi, \vspace{1em}

   Thank you so much for the information! We're begun our preliminary analysis and were hoping you could provide some additional information when you have a chance:

    \begin{itemize}
        \item A list of current bus stop locations and whether you anticipate changing any of them
        \item Any fixed stops where buses are required to stop every time
        \item Addresses of students who require bus monitors
        \item Details on the bus fleet: types of buses, capacities, ranges, and quantities of each type
        \item Information on driver scheduling:
        \begin{itemize}
            \item Does each driver typically use the same bus all day?
            \item Are there scheduled driver/bus changes during the day?
            \item Do you have an extraboard or pool of flexible/on-call drivers?
        \end{itemize}
        \item The start times for each school
        \item Any known traffic constraints or patterns that significantly affect specific routes or times
        \item Where buses start from (eg. bus yard locations, or whether they're parked at the schools themselves)
    \end{itemize}

    Please let us know if it’s possible to gather this information, or if any of these items need clarification. We’re looking forward to working with you! \vspace{1em}

    Best, \\
    Diego Temkin
\end{emailbox}

\begin{emailbox}{Thank you + Next Steps}
    \textbf{From:} Lorena Tovar \\
    \textbf{To:} Diego Temkin \\
    \textbf{Cc:} Lincoln Lynch; Riccardo Fiorista; Robert Tremblay; John P Attanuchi \\
    \textbf{Subject:} Re: Thank you + Next Steps

    \vspace{1em}
    Hi All, 
    
    Is there a way to include the \acrlong{mkv} students in the analysis?  For context, \acrlong{mkv} is a federal law that ensures students experiencing homelessness have immediate access to school and educational stability, so when kids in this category find housing outside of the district, this law ensures they can remain in their classroom until the end of the school year, which is why farther locations were on the original sheet.  We currently have about \textit{[quantity redacted]} \acrshort{mkv} students in the district. \vspace{1em}
    
    Thanks! \\
    Lorena
\end{emailbox}

\begin{emailbox}{Thank you + Next Steps}
    \textbf{From:} Lincoln Lynch \\
    \textbf{To:} Lorena Tovar \\
    \textbf{Cc:} Diego Temkin; Riccardo Fiorista; Robert Tremblay; John P Attanuchi \\
    \textbf{Subject:} Re: Thank you + Next Steps

    \vspace{1em}
    \textit{[paraphrased: Lincoln noted that this analysis focuses on big-bus service, that \gls{mkv} students who remain in-district are transported by minivan, and that the district does not transport \gls{mkv} students residing outside the city limits via big yellow buses.]}
    
    \noindent\rule{\textwidth}{1pt}

    Lincoln Lynch IV, Executive Director of Finance and Operations \\
    \Acrlong{fps} - Business Office \\
    \textit{[district office address redacted]} \\
    \textit{[email redacted]} \vspace{1em}

    \textit{[\gls{fps} standard EEO footer; reproduced once here, omitted from subsequent emails.]}
\end{emailbox}

\begin{emailbox}{Thank you + Next Steps}
    \textbf{From:} Lorena Tovar \\
    \textbf{To:} Lincoln Lynch \\
    \textbf{Cc:} Diego Temkin; Riccardo Fiorista; Robert Tremblay; John P Attanuchi \\
    \textbf{Subject:} Re: Thank you + Next Steps

    \vspace{1em}
    I think I am thinking more about the users than the buses themselves, but this makes sense-- thanks Lincoln. \vspace{1em}

    @Diego-- I can help with the start/end times for each school. I'll send by this evening. 
\end{emailbox}

\begin{emailbox}{Thank you + Next Steps}
    \textbf{From:} Lincoln Lynch \\
    \textbf{To:} Lorena Tovar \\
    \textbf{Cc:} Diego Temkin; Riccardo Fiorista; Robert Tremblay; John P Attanuchi \\
    \textbf{Subject:} Re: Thank you + Next Steps \\
    \textbf{Attached:} \textit{[attachment name redacted]}

    \vspace{1em}
    Hi, good afternoon. I have been working on obtaining the information. Here are my responses below in bold:

    \begin{itemize}
        \item A list of current bus stop locations and whether you anticipate changing any of them \textbf{A list of all students and their assigned bus stop and school is attached. There are two lines for every student. One for the morning pick up stop and one for the afternoon drop off stop. Please remove student names and student ID numbers.}
        \item Any fixed stops where buses are required to stop every time \textbf{All buses are required to stop at all bus stops listed on the excel sheet that is attached}
        \item Addresses of students who require bus monitors \textbf{All buses that have a M in front of it have a monitor assigned to the bus. Let me know if this will be enough information for you.}
        \item Details on the bus fleet: types of buses, capacities, ranges, and quantities of each type \textbf{School Bus list in Excel is attached}
        \item Information on driver scheduling:
        \begin{itemize}
            \item Does each driver typically use the same bus all day? \textbf{Yes}
            \item Are there scheduled driver/bus changes during the day?\textbf{ No changes during the day}
            \item Do you have an extraboard or pool of flexible/on-call drivers? \textbf{We do have several drivers we call spare drivers that are on call.}
        \end{itemize}
        \item The start times for each school \textbf{School start times are attached.} 
        \item Any known traffic constraints or patterns that significantly affect specific routes or times \textbf{This is tough to determine. There is constant traffic throughout the City.}
        \item Where buses start from (e.g., bus yard locations, or whether they're parked at the schools themselves) \textbf{\textit{[bus yard address redacted]}}
    \end{itemize}

    \noindent\rule{\textwidth}{1pt}

    Lincoln Lynch IV, Executive Director of Finance and Operations \\
    \Acrlong{fps} - Business Office \\
    \textit{[district office address redacted]} \\
    \textit{[email redacted]} \vspace{1em}

    \textit{[\gls{fps} standard EEO footer; reproduced once here, omitted from subsequent emails.]}
\end{emailbox}

\begin{emailbox}{Thank you + Next Steps}
    \textbf{From:} Diego Temkin \\
    \textbf{To:} Lincoln Lynch; Lorena Tovar \\
    \textbf{Cc:} Riccardo Fiorista; Robert Tremblay; John P Attanuchi \\
    \textbf{Subject:} Re: Thank you + Next Steps

    \vspace{1em}
    Hi Lincoln, \vspace{1em}

    Hope you're doing well with all the snow and thank you so much for all your help throughout this! As a quick update, we've worked on analyzing the data and iterating on our routing formulation, but we have a few additional questions we're hoping you'd be able to answer:

    \begin{itemize}
        \item Across your fleet, do you prefer spreading students across your entire fleet of buses, or do you prefer using less buses and filling them up completely?
        \item Do you currently do "trip chaining" (where a single bus may pick up students from multiple schools, then stop at multiple different schools)? Is this something you would want to do?
        \item Should special precautions be taken for students that require wheelchairs (or other specific accommodation), such as ensuring that only certain buses that can accommodate these students pick them up? If so, we would need data on which students require this access.
        \item How much do you ``value'' student pickups? Or, in other words, for a certain distance that a bus travels, how much can be spent on your end (including driver's salary, wear and tear, etc.)? This doesn't have to be an exact number, but an estimate would be helpful for us to generate different scenarios and adjust them to your needs later.
        \begin{itemize}
            \item Additionally, does this number vary depending on the student (such as based on certain demographic group, legal criteria, etc.)? If it doesn't, or if this question is unquantifiable and it's impossible to say a number, that's fine too!
        \end{itemize}
        \item What output do you need for the final routing data to be in? Would \acrshort{gtfs} data work, or is there a specific format everything would need to be in?
    \end{itemize}

    There's no rush on getting back to us, but any help would be appreciated. Thanks! \vspace{1em}

    Best, \\
    Diego Temkin
\end{emailbox}

\begin{emailbox}{Thank you + Next Steps}
    \textbf{From:} Lincoln Lynch \\
    \textbf{To:} Diego Temkin \\
    \textbf{Cc:} Lorena Tovar; Riccardo Fiorista; Robert Tremblay; John P Attanuchi \\
    \textbf{Subject:} Re: Thank you + Next Steps

    \vspace{1em}
    Hi Diego. Doing well here and hopefully you all are too. Please see my responses in bold below \vspace{1em}

    Hope you're doing well with all the snow and thank you so much for all your help throughout this! As a quick update, we've worked on analyzing the data and iterating on our routing formulation, but we have a few additional questions we're hoping you'd be able to answer:

    \begin{itemize}
        \item Across your fleet, do you prefer spreading students across your entire fleet of buses, or do you prefer using less buses and filling them up completely? \textbf{We prefer spreading students out so they do not sit 2 or 3 to a seat.}
        \item Do you currently do "trip chaining" (where a single bus may pick up students from multiple schools, then stop at multiple different schools)? Is this something you would want to do? \textbf{Yes we do this and we call it a three tiered system. We have buses that transport High School students then Middle School students then Elementary School students.}
        \item Should special precautions be taken for students that require wheelchairs (or other specific accommodation), such as ensuring that only certain buses that can accommodate these students pick them up? If so, we would need data on which students require this access. \textbf{I will work on this. We only have a few buses that have wheelchair capability.}
        \item How much do you ``value'' student pickups? Or, in other words, for a certain distance that a bus travels, how much can be spent on your end (including driver's salary, wear and tear, etc.)? This doesn't have to be an exact number, but an estimate would be helpful for us to generate different scenarios and adjust them to your needs later. \textbf{Hi we do not attach a value. We look at a per day cost per bus at this point based on the lease cost we pay. \textit{[\gls{fps} reported a per-bus daily cost in the low hundreds of dollars]}}
        \begin{itemize}
            \item Additionally, does this number vary depending on the student (such as based on certain demographic group, legal criteria, etc.)? If it doesn't, or if this question is unquantifiable and it's impossible to say a number, that's fine too!
        \end{itemize}
        \item What output do you need for the final routing data to be in? Would \acrshort{gtfs} data work, or is there a specific format everything would need to be in? \textbf{I have never worked with \acrshort{gtfs} data but it looks like it contains \acrshort{csv} files? Is that correct? Ideally I would not have to install additional softwares/platforms to view the data.}
    \end{itemize}

    There's no rush on getting back to us, but any help would be appreciated. Thanks! \textbf{Hope this helps and I will work on the special precautions information for you.} \vspace{1em}

    \textbf{Sincerely}, \vspace{1em}
    
    \textbf{Lincoln}

    \noindent\rule{\textwidth}{1pt}

    Lincoln Lynch IV, Executive Director of Finance and Operations \\
    \Acrlong{fps} - Business Office \\
    \textit{[district office address redacted]} \\
    \textit{[email redacted]} \vspace{1em}

    \textit{[\gls{fps} standard EEO footer; reproduced once here, omitted from subsequent emails.]}
\end{emailbox}

\begin{emailbox}{Thank you + Next Steps}
    \textbf{From:} Diego Temkin \\
    \textbf{To:} Lincoln Lynch; Lorena Tovar \\
    \textbf{Cc:} Riccardo Fiorista; Robert Tremblay; John P Attanuchi \\
    \textbf{Subject:} Re: Thank you + Next Steps

    \vspace{1em}
    Hi Lincoln, \vspace{1em}

    I hope you had a good weekend, and thank you for your response. \vspace{1em}

    To clarify the term ``trip chaining'': I’m referring to a scenario in which a single bus picks up multiple students in the same general area and then stops at multiple different schools. This could be useful for areas where many students in a given area attend various schools further away, reducing the number of buses needed to serve them. We're also operating under the assumption that students of different levels (elementary, middle, and high school) should not be assigned to the same bus at the same time. \vspace{1em}

    Regarding data outputs, we can certainly provide \acrshort{csv} files. Please let us know if you have any preferred formats, field names, or data structures, and we will align our exports accordingly. \vspace{1em}

    As a quick update, we are in the process of finalizing the mathematical formulation for the school bus routing problem and are constructing a graph-based network representation of the City of Framingham to support our modeling. We are also working on integrating historical traffic data from Waze to help ensure that our travel-time estimates are as accurate as possible. \vspace{1em}

    Lorena mentioned that some visualizations might be useful, so I've attached several plots to this email, including:
    \begin{itemize}
        \item A map of student distances to schools
        \item A preliminary test of our routing formulation
        \item The computational network representation of Framingham that we have generated so far
    \end{itemize}

    We expect to have preliminary results available soon. To support validation and comparison with the existing system, would it be possible to obtain:
    \begin{itemize}
        \item A copy of all current bus routes,
        \item A list of all students that require monitors and/or wheelchair-accessible buses, and
        \item \acrshort{gps} traces of school buses over a certain time period?
    \end{itemize}

    Additionally, if you have a list of all district-approved bus stops (including address information), that would be extremely helpful for our analysis. I understand that you are all very busy, so please let me know if getting this information is feasible or if we should proceed differently. Thanks! \vspace{1em}

    Best, \\
    Diego Temkin
\end{emailbox}

\begin{emailbox}{Thank you + Next Steps}
    \textbf{From:} Lorena Tovar \\
    \textbf{To:} Lincoln Lynch \\
    \textbf{Cc:} Robert Tremblay; John P Attanuchi; Diego Temkin; Riccardo Fiorista \\
    \textbf{Subject:} Re: Thank you + Next Steps

    \vspace{1em}
    Hi Lincoln, \vspace{1em}

    I just wanted to gently follow up on this. John mentioned in a separate thread that, in terms of data priorities, even a small sample of \acrshort{gps} traces would be very helpful for calibrating bus stop dwell times, since they currently have little or no data available for that. \vspace{1em}

    Were you able to get in touch with the bus vendor about this? Please let me know if there’s anything I can do to help. \vspace{1em}

    Many thanks! \\
    Lorena
\end{emailbox}

\begin{emailbox}{Thank you + Next Steps}
    \textbf{From:} Lincoln Lynch \\
    \textbf{To:} Lorena Tovar \\
    \textbf{Cc:} Robert Tremblay; John P Attanuchi; Diego Temkin; Riccardo Fiorista \\
    \textbf{Subject:} Re: Thank you + Next Steps \\
    \textbf{Attached:} \textit{[attachment name redacted]}

    \vspace{1em}
    Diego and team, \vspace{1em}

    Two documents attached.

    1. The attached student assignment sheet has all bus stops in the relevant stop-location column. \vspace{1em}

    2. The second attached document is an in/out operations report. It shows each bus arrival time for both AM and PM service and includes the expected arrival time at school. The first tab includes bus-school records for a sample period, and the other tabs summarize late arrivals by school. \vspace{1em}
    
    I hope this data helps. Please let me know if there are more data points needed and I can work on it asap. Have a great weekend everyone! \vspace{1em}
    
    Lincoln

    \noindent\rule{\textwidth}{1pt}

    Lincoln Lynch IV, Executive Director of Finance and Operations \\
    \Acrlong{fps} - Business Office \\
    \textit{[district office address redacted]} \\
    \textit{[email redacted]} \vspace{1em}

    \textit{[\gls{fps} standard EEO footer; reproduced once here, omitted from subsequent emails.]}
\end{emailbox}

\begin{emailbox}{Thank you + Next Steps}
    \textbf{From:} Robert Tremblay \\
    \textbf{To:} Diego Temkin  \\
    \textbf{Cc:} Lincoln Lynch; Lorena Tovar; John P Attanuchi;  Riccardo Fiorista \\
    \textbf{Subject:} Re: Thank you + Next Steps

    \vspace{1em}
    Thank you for the update, Diego! \vspace{1em}

    When do you think we can plan to review this together as a follow-up to our last meeting? I am eager to learn about your discoveries and hopefully some efficiencies that we can consider as we start planning for the fall. \vspace{1em}

    Thanks again for your hard work on this! \vspace{1em}

    My best, \vspace{1em}

    Bob \vspace{1em}

    \noindent\rule{\textwidth}{1pt}

    \textbf{Dr. Robert A. Tremblay} \\
    Superintendent, \acrlong{fps} \\
    Office: \textit{[phone and email redacted]} \vspace{1em}
    
    \acrlong{fps} \\
    \textit{[district office address redacted]}
\end{emailbox}
}
